\naslov{Nelinearne enačbe prvega reda}

Rešujemo enačbo oblike
\begin{equation}
  \label{eq:ne}
  F(x, y, u, u_x, u_y) = 0.
\end{equation}
Tu je $F \in \zvodvi{1}$ funkcija petih realnih spremenljivk in $u = u(x, y)$
iskana funkcija.
Vpeljemo oznake $p = u_x$, $q = u_y$.
Normala na ploskev $\{z = u(x, y)\}$ je $(u_x, u_y, -1) = (p, q, -1)$.
Enačba~\eqref{eq:ne} je zveza med točko $(x, y, z)$ na $\Gamma_u$ in normalo.
Družino potencialnih normal parametriziramo z $\vec{n}(\lambda) = (p(\lambda),
q(\lambda), -1)$.
Za vsak izbrani $\lambda$ dobimo potencialno tangentno ravnino na ploskev.
Enačba te ravnine, ki jo označimo s $\Pi_{\vec{r}_0, \lambda} = \Pi_\lambda$, je
\[
  \sk{\vec{r} - \vec{r}_0, \vec{n}(\lambda)} = 0.
\]

\begin{definicija}
  Ogrinjača ravnin $\Pi_\lambda$ se imenuje \pojem{Mongeov stožec} v točki
  $(x_0, y_0, z_0)$.
\end{definicija}

Ogrinjačo dobimo tako, da enačbo ravnin odvajamo po $\lambda$ in enačimo z $0$.
Imamo torej $\sk{\vec{r} - \vec{r}_0, \vec{n}(\lambda)} = 0$ in $\sk{\vec{r} -
  \vec{r}_0, \vec{n}'(\lambda)} = 0$.
Vektor $\vec{r} - \vec{r}_0$ je torej vzporeden z vektorskih produktom
$\vec{n}(\lambda) \times \vec{n}'(\lambda)$.
Če je $F(x_0, y_0, z_0, p(\lambda), q(\lambda)) = 0$ in odvajamo po $\lambda$,
dobimo
\[
  F_p p' + F_q q' = 0.
\]
Sledi $(q', -p') \parallel (F_p, F_q)$, torej (v kombinaciji s prejšnjo enačbo)
$\vec{r} - \vec{r}_0 \parallel (F_p, F_q, p F_p + q F_q)$.
Dobili smo parametrizacijo Mongeovega stožca
\begin{align*}
  \vec{r} &= \vec{r}_0 + \mu (P(\lambda), Q(\lambda), R(\lambda)), \\
  P(\lambda) &= F_p(x_0, y_0, z_0, p(\lambda), q(\lambda)), \\
  Q(\lambda) &= F_q(x_0, y_0, z_0, p(\lambda), q(\lambda)), \\
  R(\lambda) &= p(\lambda) P(\lambda) + q(\lambda) Q(\lambda).
\end{align*}
Privzemimo, da imamo integralno ploskev $S = \{z = u(x, y)\}$.
Naj bo $\gamma(t) = (x(t), y(t), z(t))$ neka krivulja na tej ploskvi.
Tedaj $\dot{\gamma}(t)$ leži na tangentni ravnini, zato $\vec{r}_0 +
\dot{\gamma}(t)$ leži v Mongeovem stožcu.
Posledično mora veljati (oz.~lahko zahtevamo)
\begin{align*}
  \dot{x} &= F_p, \\
  \dot{y} &= F_q, \\
  \dot{z} &= p F_p + q F_q,
\end{align*}
kjer je $p = p(\lambda)$ in $q = q(\lambda)$ za neki (ne vemo kateri) $\lambda$.
Ker $p$ in $q$ v splošnem ne poznamo, ju obravnavamo kot novi neznanki.
Vsaki točki na $\gamma$ bomo dodali še en košček ravnine, določene z normalo
$(p, q, -1)$, s čimer dobimo \pojem{karakteristični trak}.

\begin{trditev}
  Če krivulja $\gamma(t) = (x, y, z)$ zadošča pogoju $\dot{z} = p \dot{x} + q
  \dot{y}$, leži v karakterističnem traku, določenem s $p = u_x$ in $q = u_y$.
\end{trditev}

\begin{proof}
  Računamo $\sk{(\dot{x}, \dot{y}, \dot{z}), (p, q, -1)} = \sk{(\dot{x}, \dot{y},
  p \dot{x} + q \dot{y}), (p, q, -1)} = 0$.
\end{proof}


% LocalWords:  Mongeov Mongeovega Mongeovem
